\section{Conclusion}

We ask whether small models “understand’’ financial news in the limited, testable sense of aligning sentiment with same-day stock moves. Our evidence suggests that a compact, finance-tuned model (FinBERT) does extract a weak but economically meaningful signal, whereas a similarly compact general model (DistilBERT) does not. For many practical settings, \emph{domain adaptation beats raw size}.

\paragraph{Limitations.}
We analyze headlines, not full articles; we consider same-day open$\to$close returns only; and we do not control for confounders (macro releases, earnings timing). Correlations are small and should not be read as trading advice. Re-running the pipeline on different date windows will change sample sizes and noise characteristics.

\paragraph{Future work.}
Extend to multi-day horizons, intraday bars, richer event controls, and more compact domain models (e.g., finance-tuned DistilBERT or MiniLM). Explore cross-sectional predictability and interaction with volatility regimes.

\paragraph{Ethics and reproducibility.}
We cite all software and models used \citep{wolf2020transformers,araci2019finbert,prosusfinbert,sanh2019distilbert,sst2,feedparser,stooq} and provide intermediate CSVs and figures to enable replication.
